\documentclass{chsmc}
\chsmcsetup{
  year = 2023,
  part = 1,
  type = A,
  show-answers = false,
}
\begin{document}

\section{填空题：本大题共 8 小题，每小题 8 分，共 64 分.}

\begin{enumerate}
  \item 设复数 $z = 9 + 10\ii$（$\ii$ 为虚数单位），若正整数 $n$ 满足
    $|z^n| \leqslant 2023$，则 $n$ 的最大值为 \blankline
  \item 若正实数 $a$, $b$ 满足 $a^{\lg b} = 2$，$a^{\lg a}\cdot b^{\lg b} = 5$.
    则 $(ab)^{\lg ab}$ 的值为 \blankline
  \item 将一枚均匀的骰子独立投掷三次，所得的点数依次记为 $x$、$y$、$z$，
    则事件“$\mathrm{C}_7^x < \mathrm{C}_7^y < \mathrm{C}_7^z$”发生的概率为 \blankline
  \item 若平面上非零向量 $\vec{\alpha}$, $\vec{\beta}$, $\vec{\gamma}$ 满足
    $\vec{\alpha} \perp \vec{\beta}$，
    $\vec{\beta}\cdot\vec{\gamma} = 2|\vec{\alpha}|$，
    $\vec{\gamma}\cdot\vec{\alpha} = 3|\vec{\beta}|$，
    则 $|\vec{\gamma}|$ 的最小值为 \blankline
  \item 方程 $\sin x = \cos 2x$ 的最小的 $20$ 个正实数解之和为 \blankline
  \item 设 $a, b, c$ 为正数，$a < b$，若 $a$、$b$ 为一元二次方程
    $ax^2 - bx + c = 0$ 的两个根，且 $a$、$b$、$c$ 是一个三角形的三边长，
    则 $a + b - c$ 的取值范围是 \blankline
  \item 平面直角坐标系 $xOy$ 中，已知圆 $\Omega$ 与 $x$ 轴、$y$ 轴均相切，
    圆心在椭圆 $\Gamma\colon \dfrac{x^2}{a^2} + \dfrac{y^2}{b^2} = 1$
    $(a > b > 0)$ 内，且 $\Omega$ 与 $\Gamma$ 有唯一的公共点 $(8,9)$.
    则 $\Gamma$ 的焦距为 \blankline
  \item 八张标有 A、B、C、D、E、F、G、H 的正方形卡片构成下图，
    现逐一取走这些卡片，要求每次取走一张卡片时，
    该卡片与剩下的卡片中至多有一张有公共边
    （例如可按 D、A、B、E、C、F、G、H 的次序取走卡片，
    但不可按 D、B、A、E、C、F、G、H 的次序取走卡片），
    则的不同次序的取走这八张卡片的不同次序的数目为 \blankline
\end{enumerate}

\section{解答题：本大题共 3 个小题，满分 56 分. 解答应写出文字说明、证明过程或演算步骤.}

\begin{enumerate}[resume]
  \item (本题满分 16 分) 平面直角坐标系 $xOy$ 中，抛物线
    $\Gamma\colon y^2 = 4x$，$F$ 为 $\Gamma$ 的焦点，$A$、$B$ 为 $\Gamma$
    上的两个不重合的动点，使得线段 $AB$ 的一个三等分点 $P$ 位于线段 $OF$ 上（含端点），
    记 $Q$ 为线段 $AB$ 的另一个三等分点. 求点 $Q$ 的轨迹方程.
  \item (本题满分 20 分) 已知三棱柱 $\Omega\colon ABC\text{-}A_1B_1C_1$
    的 $9$ 条棱长均相等，记底面 $ABC$ 所在平面为 $\alpha$.
    若 $\Omega$ 的另外四个面（即面 $A_1B_1C_1$, $ABB_1A_1$, $ACC_1A_1$,
    $BCC_1B_1$）在 $\alpha$ 上投影面积从小到大重排后依次为
    $2\sqrt{3}$, $3\sqrt{3}$, $4\sqrt{3}$, $5\sqrt{3}$.
    求 $\Omega$ 的体积.
  \item (本题满分 20 分) 求出所有满足下面要求的不小于 $1$ 的实数 $t$：
    对任意 $a, b \in [-1,t]$，总存在 $c, d \in [-1,t]$，
    使得 $(a+c)(b+d) = 1$.
\end{enumerate}

\end{document}